% !PRINTED 673
\weberpage{673}
Diese einfache Bemerkung liefert uns den folgenden wichtigen Satz:

2. \gesperrt{Es giebt in einem algebraischen K\"orper $\Omega$ nur eine endliche Anzahl ganzer Zahlen $\eta$ von der Beschaffenheit, dass die absoluten Werthe der conjugirten Zahlen $\eta_s$ unter einer gegebenen Grenze liegen.}

Wir machen jetzt von dem allgemeinen Satze \S.\,155, (25) Gebrauch, nach dem, wenn $f(x_1,x_2,\dots,x_n)$ irgend eine \gesperrt{positive quadratische Form} von $n$ Variablen mit der Determinante $D$ ist, die Variablen $x_i$ sich als ganze rationale Zahlen, nicht alle verschwindend, so bestimmen lassen, dass
\[
f(x_1,x_2,\dots,x_n) < n\sqrt[n]{0{,}404\,D},
\]
d.\,h. kleiner als eine von der Determinante $D$ allein abh\"angige Zahl wird.

Diesen Satz wenden wir, wie im \S.\,156, auf die Form an:
\[
f(x_1,x_2,\dots,x_n) = \frac{|\eta_1|^2}{c_1^2} + \frac{|\eta_2|^2}{c_2^2} + \cdots + \frac{|\eta_n|^2}{c_n^2},
\]
worin die $\eta_s$ die Linearformen (4) sind, und im Falle eines reellen $\eta_s$
\[
|\eta_s|^2 = (\eta_s)^2,
\]
im Falle eines imagin\"aren Paares $\eta_s, \eta_s'$
\[
|\eta_s|^2 = |\eta_s'|^2 = \eta_s\eta_s'.
\]
Im letzteren Falle ist $\eta_s\eta_s'$ die Summe zweier reeller Quadrate, und demnach ist, wenn die $c_s$ reell angenommen werden, $f$ eine positive Form. Wir wollen \"uber die bis jetzt willk\"urlichen reellen Gr\"ossen $c_s$ noch festsetzen, dass, wenn $\eta_s, \eta_s'$ ein conjugirt imagin\"ares Paar bilden,
\weq{5}{c_s = c_s'}
sein soll; ausserdem wollen wir die $c_s$ noch der Bedingung unterwerfen:
\weq{6}{c_1 c_2 \dots c_n = 1.}
Die Determinante der Function $f$ bilden wir, wie im \S.\,156, dadurch, dass wir sie zun\"achst als Form der Variablen $\eta_1, \eta_2, \dots, \eta_n$ auffassen, und die Determinante dieser Form ist wegen (6) gleich $\pm 1$. Die Determinante von $f$ in Bezug auf die Variablen $x_i$ ist also (Bd.\ I, \S.\,56), vom Vorzeichen abgesehen, gleich dem
%CONT>>
